% 公共接口与加载顺序测试的正文。由 tests/if-*.tex 提供主题与加载顺序。
% 覆盖上游公开接口、Minimalist 扩展接口与标准附录。
\providecommand{\hitdemo}{minimalist}
\providecommand{\hitorder}{after}

\ifdefined\hitmode\else\def\hitmode{theme}\fi
\def\hit@loadtheme{%
  \ifnum\pdfstrcmp{\hitdemo}{Madrid}=0
    \usetheme{Madrid}%
  \else\ifnum\pdfstrcmp{\hitdemo}{classic}=0
    \usetheme[classic]{hit}%
  \else\ifnum\pdfstrcmp{\hitdemo}{touying}=0
    \usetheme{hit}%
  \else\ifnum\pdfstrcmp{\hitdemo}{hit}=0
    \usetheme{hit}%
  \else\ifnum\pdfstrcmp{\hitdemo}{minimalist}=0
    \usetheme[minimalist]{hit}%
  \else\ifnum\pdfstrcmp{\hitdemo}{hitminimalist}=0
    \usetheme[minimalist]{hit}%
  \else\ifnum\pdfstrcmp{\hitdemo}{hitacademic}=0
    \usetheme[minimalist]{hit}%
  \else
    \usetheme[\hitdemo]{hit}%
  \fi\fi\fi\fi\fi\fi\fi
}
\ifnum\pdfstrcmp{\hitorder}{before}=0
  \usepackage{beamercmdhit}
  \hit@loadtheme
\else
  % 主题自行加载命令层，之后再次声明不得产生重复定义。
  \hit@loadtheme
  \usepackage{beamercmdhit}
\fi

\title[接口测试]{公共接口与加载顺序测试}
\englishtitle{Public Interface Test}
\titlecontext{接口矩阵}
\subtitle{副标题内容不得丢失}
\author{汇报人：待填写}
\institute{哈尔滨工业大学}
\date{2026-10-01}

% 上游公开接口：目录标题可被 \renewcommand 改写。
\renewcommand{\hitoutlinetitle}{目录（自定义标题）}

\begin{document}

\TitleSlide
\BlueTitleSlide
\OutlineSlide

\section{公开命令}

\begin{frame}{公开命令的兼容调用}
  \begin{itemize}
    \item \texttt{\textbackslash TitleSlide} 与 \texttt{\textbackslash BlueTitleSlide}
    \item \texttt{\textbackslash OutlineSlide}：无参数调用
    \item \texttt{\textbackslash hitfooter}、\texttt{\textbackslash hitvipath}
  \end{itemize}
  \hitfooter{页脚文字}
\end{frame}

\begin{frame}{Minimalist 扩展接口}
  \subhead{小标题}
  \insight{结论区在三种主题下都要能排版。}
  \notebox{注记保留 2 pt 左侧竖线。}
  \captiontext{图注使用主题淡色。}
  \begin{columns}[T,onlytextwidth]
    \begin{column}{.30\textwidth}
      \metric{\linewidth}{2.4×}{解析模型给出的加速比}
    \end{column}
    \begin{column}{.62\textwidth}
      \begin{denseitems}
        \item 字号辅助命令可用。
        \item 条目间距由 denseitems 收紧。
      \end{denseitems}
    \end{column}
  \end{columns}
\end{frame}

\FocusSlide{单参数重点页}

\section{附录与结束}

\begin{frame}{标准附录入口}
  附录使用标准 \texttt{\textbackslash appendix}，页码分母保持为正文页数。
\end{frame}

\appendix

\begin{frame}{附录内容}
  附录页的页脚与页码按主题处理。
\end{frame}

\EndSlide{感谢聆听}

\end{document}
